\section{Randomized equality at four queries with additive costs}
\label{sec:four-additive}
The fee class changes the randomized boundary. Here we prove an exclusion
at four queries for every nonnegative additive price vector. This is
stronger than agreement on a finite set of prices. It does not extend
the exclusion to arbitrary partial Boolean promises or to all submodular
final-set fees.

For $\mathcal C\in\{\mathcal A,\mathcal P\}$, define
\[
 R_{\mathcal C}(W;c)=\min_{\mu}\max_{w\in W}
   \sum_T\mu_T\sum_{q\in P_T(w)}c_q,\qquad c\in\mathbb R_{\geq0}^{Q}.
\]
The private seed is independent of $w$. For $\mathcal P$ each seed may
choose its own global permutation; its two answer branches share that
permutation and may permanently skip queries using the observations so
far. A lower-bound prior used in a certificate is not a population
assumption. The maximum is over the original correctness promise, even
at inputs receiving zero probability in that prior.

\begin{theorem}[Four-query additive equality, computer-assisted]
\label{thm:four-additive}
Let $W_0,W_1$ be the nonempty, disjoint projections of two complete
partition fibres with an unrestricted number of classes, on at most
four indexed pair-equality queries. For every finite nonnegative real
additive price vector $c$,
\[
 R_{\mathcal A}(W;c)=R_{\mathcal P}(W;c).
\]
Hidden records, shared endpoints, zero prices and price ties are allowed.
The statement concerns worst-input expected actual purchased cost at
zero error. It does not assert equality under every specified input prior.
\end{theorem}

\begin{proof}
We give the reduction, full-policy premise and continuous-price
certificate obligations. The accompanying packet supplies the finite
assignments and complete analytic case arguments. Its implementation
checks the stated obligations using exact arithmetic; the analytic
arguments are not replaced by a formal proof kernel.

\emph{Complete source cover.}
Collapse all prior SAME components, including paths through hidden
records. The eight marked endpoint occurrences induce a partition $\pi$.
For an answer word $w\in\{0,1\}^4$, merge the positive query edges.
The word is possible exactly when every zero-answer edge and every
prior DIFFERENT pair still has separate ends. Necessity follows from
transitivity. For sufficiency assign each remaining component a fresh
class, and assign fresh classes to hidden-only components. The number
of classes is unrestricted, so there is no further extension condition.

For a fixed $\pi$, let $B_\pi$ be its answer support without forbidden
pairs. A forbidden pair between endpoint-containing blocks defines a
filter $H_j\subseteq B_\pi$. Starting with $\{B_\pi\}$, the iteration
\[
 \mathcal L_j=\mathcal L_{j-1}\cup
       \{S\cap H_j:S\in\mathcal L_{j-1}\}
\]
covers every finite forbidden-pair conjunction, including redundant
filters. All $4140$ restricted-growth strings on eight positions give
$465$ distinct nonempty supports. Every disjoint pair is realized on
the same eight physical records with its two label-specific priors.
There are $42864$ ordered tasks, partitioned into $1356$ orbits under
simultaneous query-coordinate permutations and label exchange.
Answer-bit complementation is not a physical symmetry used in this cover.

\emph{Full policy cover.}
Repeated, publicly fixed-answer and post-pure queries can be erased
without increasing a nonnegative additive cost. In $\mathcal P$ erasure
is a legal permanent skip. At a mixed feasible state, the adaptive
recursion chooses any remaining informative query and combines both
answer continuations. For a fixed permutation, the priority recursion
unites SKIP on the same feasible set with TAKE on both answer children,
using the same remaining suffix on both branches. Pure states stop.
Taking the union over all $24$ orders covers the full priority class.
The finite union of input-specific null sets of erroneous seeds can
be removed from any zero-error randomized policy. Seedwise
normalization therefore reduces it to a distribution on these finite
families. A fresh recursion reconstructs all $39113$ adaptive and
$33781$ priority paid-set vectors. All $629694$ original-input paths
are checked for correctness, purchased sets, no repeated queries and
the shared global order. No chosen-order or restricted-input baseline
is substituted for $\mathcal P$.

\emph{The majority of sources.}
For any normalized monotone fee $F$, put
\[
 K_F=\max_{w\in W}\min_{S\text{ certifies }w} F(S).
\]
Every correct path is a certificate, so $K_F\leq R_{\mathcal A}(W;F)$.
The set $I_F=\{S:F(S)\leq K_F\}$ is a nonempty downset in $2^{[4]}$
and contains a certificate for every input. There are exactly $167$
nonempty downsets. For $1300$ orbits, each relevant downset admits a
legal original priority policy whose path sets all belong to it.
Checking all $155995$ relevant source/downset cells yields
$R_{\mathcal P}\leq K_F\leq R_{\mathcal A}$. For another $31$ orbits,
each complete adaptive vector has a legal priority vector contained
in it on every original input. Replacing each seed's policy proves
the same inequality. The reverse inequality follows throughout from
$\mathcal P\subseteq\mathcal A$.

\emph{Structural and additive-only sources.}
Table~\ref{tab:four-proof-cover} gives the disjoint remainder. For the
$14$ structural orbits, complete two-column, three-column and chain
saddle arguments, or common-input retractions, imply equality for all
normalized monotone fees. The packet includes their complete ordered
case splits, generator signs and zero-boundary arguments. Polynomial
identities and nonnegative-product decompositions are checked for
$64$, $105$ and $52$ sign obligations in the three scalar saddle
lemmas, respectively. The physical source bindings are separately
checked on the full original families ($2336$ input obligations).
Auxiliary response encodings in a structural template are individually
bound between its lower core and upper envelope; they are not added
to the physical orbit group. The two cost-preorder regions use opposite
scalar inequalities and include their shared boundary.

For clarity, the remaining $11$ orbits require only additive costs.
At a fixed nonnegative price vector, a sufficient certificate is a
probability vector $\lambda$ on original inputs and a probability
vector $p$ on legal priority columns with common value $L$, satisfying
\begin{align}
 \sum_w\lambda_w\,c(P_A(w))&\geq L
       &&\text{for every full adaptive column }A,\label{eq:four-dual}\\
 \sum_jp_j\,c(P_{P_j}(w))&\leq L
       &&\text{for every original input }w.\label{eq:four-upper}
\end{align}
These imply $R_{\mathcal A}\geq L\geq R_{\mathcal P}$ directly.
For constant rational $\lambda,p$, the inequalities are homogeneous
linear forms in $c$. On a closed price-order cone
$c_{\pi_0}\leq\cdots\leq c_{\pi_3}$, write
\[
 c_{\pi_j}=\sum_{k=0}^{j}t_k,\qquad t_k\geq0.
\]
Nonnegative coefficients in the $t$ variables prove an inequality
on the entire cone, with ties and zero prices included.

The retained-library orbit uses all $29$ unchanged rational
prior/priority-mixture certificates. Every one of the $24$ order cones
lies in a certified region, as checked on its four suffix-indicator
generators. The exact normals are reconstructed from all original
adaptive columns and all original-input priority upper inequalities;
the library is not justified by its old sampled prices.

Nine more orbits use $24$ complete closed order cones each. Each cone
has a pure saddle certificate or a sequence of row and column
dominance inequalities with nonnegative coefficients. Deleting a
column replaces it by a no-more-expensive active column on every
active row. Deleting a row maps it to a no-cheaper active row
simultaneously for all active columns. Following the row chains gives
one common map for every mixture of terminal columns. Each terminal
column then lifts to an original legal priority column on all hard
rows. This preserves the full worst-input randomized value. In the
one orbit with a preliminary deleted input, its paid set is contained
in another input's paid set for \emph{every} original adaptive and
priority column. Its objective constraint is therefore redundant for
both classes, while it remains a correctness obligation. All $16634$
dominance steps, $556560$ scalar signs and $3008$ pure-saddle scalar
signs are replayed from the original paid sets.

\begin{table}[ht]
\centering\small
\begin{tabular}{@{}p{0.67\linewidth}r@{}}\toprule
Proof obligation & Source orbits\\\midrule
Certificate-threshold downsets & $1300$\\
Same-input full-policy subset replacement & $31$\\
Two-column structural saddle & $2$\\
Three-column saddle and signed template bindings & $1+4$\\
Common retraction; two preorder regions; AND/hint; prefix & $1+1+1+1$\\
Chain structural saddle & $3$\\
Retained rational library, all $29$ regions & $1$\\
Pure saddles and common-row dominance, $24$ cones each & $1+8$\\
Balanced two-input certificate and $22$ old closed cones & $1$\\\midrule
Total, with disjoint memberships & $1356$\\\bottomrule
\end{tabular}
\caption{Complete four-query cover. The first six rows prove equality
for the full monotone fee class on their stated sources. The last
three rows complete the nonnegative additive theorem only. Each orbit
is bound to its full original policy families.}
\label{tab:four-proof-cover}
\end{table}

\emph{The final source and its boundary faces.}
To make the remaining rational dependence explicit, one orbit has
negative words $\{0,3,4,7,8,11\}$ and positive words $\{1,15\}$,
where $w=\sum_{q=0}^3 2^qw_q$. Write its prices as $(a,b,c,d)$.
Twenty-two order cones are closed by the preceding dominance method.
The two remaining orders are interchanged by swapping $c,d$; this
involution preserves the source and both complete policy families.
It suffices to treat $0\leq a\leq b\leq c\leq d$.
Two legal original priority policies, denoted $U,V$, have paid masks
$(3,14)$ and $(7,12)$ respectively on positive inputs $(1,15)$.
Their full original-input traces and global orders are checked.
Put
\[
 D=b+c,\quad S=a+b-d,\quad
 N=b(a+b+c)+c(c+d).
\]
On $S\geq0$ and $D>0$, use weights $S/D$ and $(c+d-a)/D$
on $U,V$, and lower prior $b/D,c/D$ on inputs $1,15$.
All weights are nonnegative and sum to one. Both positive inputs
have upper value $N/D$. The common row map
\[
 \tau(0)=\tau(1)=\tau(8)=1,\qquad
 \tau(3)=\tau(4)=\tau(7)=\tau(11)=\tau(15)=15
\]
dominates every original input for both upper policies on this cone.

For the lower bound, every correct tree must buy queries $a,b$ on
input $1$ and $c,d$ on input $15$, using the opposite-label
single-coordinate witnesses $0,3$ and $11,7$ respectively. These
paths share their prefix up to the first differing queried answer.
If that first differing query is $b$, their costs are at least
$(a+b,b+c+d)$; if it is $c$, they are at least $(a+b+c,c+d)$.
The prior-weighted numerators in both cases are $N$. If it is $d$,
the latter lower numerator increases by $b(d-c)\geq0$. Extra purchases
cannot reduce cost. Thus \eqref{eq:four-dual} holds for every correct
adaptive tree, not just the retained finite columns.

On $S\leq0$, pure policy $V$ has worst cost $c+d$; the mandatory
$c,d$ purchases on input $15$ give the matching lower bound. The
packet checks all original-input upper inequalities. Closed
nonnegative parameterizations for the three regions are
\begin{align*}
 S\geq0 &: (S,b-a,c-b,d-c),\\
 S\leq0,\ c\leq a+b &: (a+b-c,b-a,c-b,d-a-b),\\
 S\leq0,\ c\geq a+b &: (a,b-a,c-a-b,d-c).
\end{align*}
Their invertible linear maps and all numerator coefficients are
checked exactly ($1576$ scalar coefficients, including zeros).
In the first region $D$ has coefficient vector $(2,2,3,2)$, so
$D=0$ is only the all-zero price vector. Use legal pure $V$ there
at cost zero, without division. On the switching face,
$N-(c+d)D=bS=0$, and the mixed weight on $U$ vanishes. This
includes the shared boundaries and closes both remaining orders.

All $1356$ orbits are now covered for the whole nonnegative real
additive orthant. Pulling prices back along a physical coordinate
permutation preserves the classes and costs. For fewer than four
queries, append fresh known-SAME zero-price dummy pairs to both
fibres. Their insertion and erasure preserve the two policy classes
and all actual costs, proving the stated at-most-four scope.
\end{proof}

\begin{corollary}[Randomized additive first-occurrence interval]
\label{cor:additive-first-interval}
Let $m_{\rm add}$ and $m_{\rm pos}$ be the least menu sizes with
$R_{\mathcal A}<R_{\mathcal P}$ in this same complete-source model,
for nonnegative additive and strictly positive additive prices,
respectively. Then
\[
 5\leq m_{\rm add}\leq m_{\rm pos}\leq7.
\]
The exact minima, and the corresponding submodular minimum, are
not determined here.
\end{corollary}
\begin{proof}
Theorem~\ref{thm:four-additive} excludes all smaller menus. On the
complete seven-query source in Section~\ref{sec:seven}, use the old
prices $(2,2,3,1/3,1,1,1)$. A correct adaptive tree has worst cost
$13/3$. The original priority mixture has upper $227/51$ and the
fixed prior supported on world indices $(0,4,5,13,15)$, with
weights $(9,1,1,1,5)/17$, gives a matching lower bound over all
$5040$ global orders with causal skips. Worlds with prior zero
are still included in every correctness/purity test. Therefore
\[
 R_{\mathcal P}-R_{\mathcal A}\geq
  \frac{227}{51}-\frac{13}{3}=\frac{2}{17}>0.
\]
The retained exact curve in Appendix~\ref{app:seven-private-curve}
also supplies the adaptive optimum. Only its legal adaptive upper
and matched full-priority value are needed for this corollary;
the upper alone is not a new proof of an exact adaptive optimum.
\end{proof}

